\section{Results}
\label{sec:results}

\subsection{The label--essay-score association moves with the definition, in direction only at some levels}
\label{sec:res-score}

A study that reports how a pause or stall rate relates to the quality of a product usually reports one definition. We
asked whether that association depends on the definition, and measured, for each definition, the Spearman
correlation between the KLiCKe writers' rates and their holistic scores, with a resampled interval
(Figure~\ref{fig:score}, Table~\ref{tab:pause}).

\paragraph{Stalls.} Under the baseline the correlation was $-0.063$ [$-0.091$, $-0.034$]. Two criteria had been set
before the association was seen. The presence of an association would count as moving if the share of definitions
with an interval excluding zero lay between 25\% and 75\%; it was 82 of 89 (92\%), so by that criterion presence did not
move. The direction would count as changing if one definition had an interval on the other side of zero; 10 did (also
when the point value is required to lie there), so the direction changed. All 10 had the shortest window,
$W = 15$~s, the level we set at half the default, and a writer-quantile $\theta$; among the 59 definitions with
other windows, 57 were negative and none was positive, and on the grid of default and precedented levels
(Table~\ref{tab:grid}) none was positive either. Across the 89 definitions the
correlation ranged from $-0.25$ to $+0.14$ (median $-0.14$), with intervals negative for 72, positive for 10 and
including zero for 7. By point value, 58 definitions gave a negative correlation of at least 0.1 in absolute value and 3
a positive one; all 36 fixed-second definitions were negative. The positive corner is not a floor effect of rare
stalls (S3).

\paragraph{Pauses.} For the second detector the reading set in advance was that the direction changes if at least one
definition has an interval above zero and one below. With fixed thresholds, the interval was below zero for 7 of the 18
definitions, all but one at 10~s, and above zero for 3, on character intervals at 2 and 5~s (Table~\ref{tab:pause}).
The direction therefore changed, but between mostly negligible values: 13 of the 18 intervals lay within $\pm 0.1$, and only two definitions at
10~s had a point value at or below $-0.10$ ($-0.101$ and $-0.127$). By the criterion for presence used for stalls, which was not
set in advance for pauses, 10 of 18 intervals (56\%) excluded zero, so presence moved as well. The 10-s values are
smaller in magnitude than the $-0.15$ to $-0.43$ reported for pause durations and latencies, a different measure
(Section~\ref{sec:related}). The 18 definitions with a writer-quantile threshold, a control for the rate of key presses
(Section~\ref{sec:grid}), correlated with the score at $+0.58$ to $+0.67$, close to the 0.62 of the rate of key presses
itself; the reading stated for all 36 definitions is in S1.

A reader who took one definition would report a small negative relation (any fixed stall threshold), a
negligible one (the stall baseline, most fixed pause thresholds), or a positive one (a 15-s stall window with a writer
quantile); a writer-quantile pause threshold would report the rate of key presses.

\begin{figure}[t]
\centering
\begin{tikzpicture}
\begin{axis}[name=raw, width=\textwidth, height=4.8cm, xmin=0, xmax=91, ymin=-0.33, ymax=0.28,
  ylabel={Spearman $\rho$}, xticklabels={}, ytick={-0.3,-0.2,-0.1,0,0.1,0.2}, grid=major, grid style={figband},
  title={(a) Stall rate and holistic score, 90 definitions ordered by $\rho$}, title style={font=\footnotesize},
  tick label style={font=\footnotesize}, label style={font=\footnotesize},
  legend style={font=\scriptsize, at={(0.02,0.98)}, anchor=north west}, legend cell align=left]
  \addplot[only marks, mark=*, mark size=1.1pt, color=figc1, error bars/.cd, y dir=both, y explicit,
    error bar style={color=figc1!60}] table[x=i, y=rho, y error plus expr=\thisrow{hi}-\thisrow{rho},
    y error minus expr=\thisrow{rho}-\thisrow{lo}, discard if not={fam}{1}] {figdata/score.csv};
  \addlegendentry{writer-quantile $\theta$}
  \addplot[only marks, mark=square*, mark size=1.1pt, color=figc2, error bars/.cd, y dir=both, y explicit,
    error bar style={color=figc2!60}] table[x=i, y=rho, y error plus expr=\thisrow{hi}-\thisrow{rho},
    y error minus expr=\thisrow{rho}-\thisrow{lo}, discard if not={fam}{0}] {figdata/score.csv};
  \addlegendentry{fixed-second $\theta$}
  \addplot[only marks, mark=star, mark size=3.5pt, thick, color=black] table[x=i, y=rho, discard if not={b0}{1}]
    {figdata/score.csv};
  \addlegendentry{baseline}
  \addplot[only marks, mark=triangle, mark size=2.2pt, color=black] table[x=i, y expr=\thisrow{hi}+0.03,
    discard if not={W15}{1}] {figdata/score.csv};
  \addlegendentry{$W = 15$~s}
  \draw[black, thin] (axis cs:0,0) -- (axis cs:91,0);
\end{axis}
\begin{axis}[name=spd, at={(raw.south)}, anchor=north, yshift=-0.9cm, width=\textwidth, height=4cm, xmin=0, xmax=91,
  ymin=-0.33, ymax=0.2, ylabel={partial $\rho$}, xticklabels={}, ytick={-0.3,-0.2,-0.1,0,0.1}, grid=major,
  grid style={figband}, title={(b) Given typing speed (median interval $< 2$~s; exploratory)},
  title style={font=\footnotesize}, tick label style={font=\footnotesize}, label style={font=\footnotesize}]
  \addplot[only marks, mark=*, mark size=1.1pt, color=figc1, error bars/.cd, y dir=both, y explicit,
    error bar style={color=figc1!60}] table[x=i, y=adj, y error plus expr=\thisrow{ahi}-\thisrow{adj},
    y error minus expr=\thisrow{adj}-\thisrow{alo}, discard if not={fam}{1}] {figdata/score.csv};
  \addplot[only marks, mark=square*, mark size=1.1pt, color=figc2, error bars/.cd, y dir=both, y explicit,
    error bar style={color=figc2!60}] table[x=i, y=adj, y error plus expr=\thisrow{ahi}-\thisrow{adj},
    y error minus expr=\thisrow{adj}-\thisrow{alo}, discard if not={fam}{0}] {figdata/score.csv};
  \addplot[only marks, mark=star, mark size=3.5pt, thick, color=black] table[x=i, y=adj, discard if not={b0}{1}]
    {figdata/score.csv};
  \draw[black, thin] (axis cs:0,0) -- (axis cs:91,0);
\end{axis}
\begin{axis}[at={(spd.south)}, anchor=north, yshift=-0.9cm, width=\textwidth, height=4cm, xmin=0, xmax=91,
  ymin=-0.33, ymax=0.2, ylabel={partial $\rho$}, xlabel={definitions (same order as panel a)}, xticklabels={},
  ytick={-0.3,-0.2,-0.1,0,0.1}, grid=major, grid style={figband},
  title={(c) Given essay length and length of observation}, title style={font=\footnotesize},
  tick label style={font=\footnotesize}, label style={font=\footnotesize}]
  \addplot[only marks, mark=*, mark size=1.1pt, color=figc1, error bars/.cd, y dir=both, y explicit,
    error bar style={color=figc1!60}] table[x=i, y=ln, y error plus expr=\thisrow{lhi}-\thisrow{ln},
    y error minus expr=\thisrow{ln}-\thisrow{llo}, discard if not={fam}{1}] {figdata/score.csv};
  \addplot[only marks, mark=square*, mark size=1.1pt, color=figc2, error bars/.cd, y dir=both, y explicit,
    error bar style={color=figc2!60}] table[x=i, y=ln, y error plus expr=\thisrow{lhi}-\thisrow{ln},
    y error minus expr=\thisrow{ln}-\thisrow{llo}, discard if not={fam}{0}] {figdata/score.csv};
  \addplot[only marks, mark=star, mark size=3.5pt, thick, color=black] table[x=i, y=ln, discard if not={b0}{1}]
    {figdata/score.csv};
  \draw[black, thin] (axis cs:0,0) -- (axis cs:91,0);
\end{axis}
\end{tikzpicture}
\caption{The stall--score association in KLiCKe (4,908 writers) for the baseline and the 89 definitions, with 95\%
intervals from resampling writers. Triangles mark the definitions with $W = 15$~s. Panels (b) and (c) keep the order of panel (a).}
\label{fig:score}
\end{figure}

\begin{table}[t]
\centering\footnotesize
\caption{The pause family: the association of the pause rate with the holistic score in KLiCKe, range over the three
segmentations for each kind of interval and threshold, and the number of the three with an interval below or above
zero. Every definition is in S3.}
\label{tab:pause}
\begin{tabular}{@{}l l r r r@{}}
\toprule
Interval & Threshold & $\rho$ & Below zero & Above zero \\
\midrule
any key & 2~s & $-0.01$ to $+0.03$ & 0 & 0 \\
 & 5~s & $-0.04$ to $0.00$ & 1 & 0 \\
 & 10~s & $-0.13$ to $-0.09$ & 3 & 0 \\
character & 2~s & $+0.01$ to $+0.05$ & 0 & 2 \\
 & 5~s & $0.00$ to $+0.03$ & 0 & 1 \\
 & 10~s & $-0.10$ to $-0.06$ & 3 & 0 \\
any key & quantile 0.95--0.995 & $+0.58$ to $+0.62$ & 0 & 9 \\
character & quantile 0.95--0.995 & $+0.62$ to $+0.67$ & 0 & 9 \\
\bottomrule
\end{tabular}
\end{table}

\subsection{Adjusting for essay length leaves no stall definition negative and most values negligible}
\label{sec:res-adjust}

Two quantities that do not contain the labels can carry the association: how much a writer writes and how fast.
Essay length correlated with the score at 0.75 [0.73, 0.76], and faster writers had higher scores (Spearman $-0.53$
between the score and the median short interval, a lower interval meaning faster typing). The length of observation
barely varied in KLiCKe (interquartile range 26--29~min) and correlated with the score at 0.10.

\paragraph{Given essay length.} Given essay length and the length of observation, the partial correlation of score and
stall rate was positive with an interval above zero for 54 definitions, included zero for 35 and was below zero for
none ($+0.02$ to $+0.13$; baseline $+0.057$ [$+0.029$, $+0.085$]); the 10 definitions that were positive without
adjustment stayed positive. By the reading set before the computation, the sign became uniform: no definition was
negative. In size, 74 of the 89 intervals lay within $\pm 0.1$, so after adjustment the stall rate carried little
association with the score. Stalls and essay length overlap by construction: under the net rule a stall is a period in
which the text does not grow, so the adjustment also removes part of what the stall rate measures. For the fixed pause
thresholds, 15 of 18 were positive and none negative, and 9 had an interval above $+0.1$.

\paragraph{Given typing speed (exploratory).} With the median of intervals under 2~s, the partial correlation of score and stall rate was
negative with an interval below zero for all 89 definitions ($-0.29$ to $-0.03$; baseline $-0.18$ [$-0.21$, $-0.15$]),
and below $-0.1$ for 71; the ten definitions with a positive unadjusted interval all had an interval below zero.
Six ways of measuring speed gave no positive interval and 2 of
534 intervals including zero (S3). For the fixed pause thresholds the speed-adjusted correlation went
the other way: positive for 11 of 18 and negative for 4.

Neither adjusted value is a causal effect (Section~\ref{sec:discussion}).

\paragraph{Reading by classes of interval.} Table~\ref{tab:read} applies the reading proposed in
Section~\ref{sec:discussion}, formed after seeing these results.
Unadjusted, the fixed stall thresholds read as below $-0.1$ and the fixed pause thresholds as negligible, while the full
stall grid depended on the definition. On the grid of default and precedented levels the reading split by the family of
$\theta$, which the grid treats as two constructs (Section~\ref{sec:grid}): the fixed 10-s threshold read as below
$-0.1$ and the writer quantile as negligible, each consistently. Given length, every stall group read as negligible, but
the fixed pause thresholds came to depend on the definition, with 9 of 18 above $+0.1$: adjustment changed the reading
as a definition does. Given speed, the full grid fell just short of a consistent reading (89.9\%).

\begin{table}[t]
\centering\scriptsize\setlength{\tabcolsep}{3pt}
\caption{Reading of the score association by classes of interval (a rule proposed after seeing the results). Each
cell gives the reading, the share of the leading class among the classified definitions, and the counts below
$-0.1$ / within $\pm 0.1$ / above $+0.1$ / imprecise. Groups with the baseline: the default and precedented grid and
its writer-quantile part; the full grid excludes it.}
\label{tab:read}
\begin{tabular}{@{}l l l l@{}}
\toprule
Group (definitions) & Unadjusted & Given length & Given speed (exploratory) \\
\midrule
Stalls, full grid (89) & depends (71\%): 48/17/2/22 & negligible (98\%): 0/74/1/14 & depends (89.9\%): 71/8/0/10 \\
Stalls, fixed $\theta$ (36) & below (100\%): 36/0/0/0 & negligible (100\%): 0/33/0/3 & below (100\%): 36/0/0/0 \\
Stalls, writer quantile (53) & depends (54\%): 12/17/2/22 & negligible (97\%): 0/41/1/11 & depends (81\%): 35/8/0/10 \\
Stalls, default and precedented (12) & depends (62\%): 3/5/0/4 & negligible (100\%): 0/9/0/3 & below (100\%): 12/0/0/0 \\
\quad of which fixed 10~s (3) & below (100\%): 3/0/0/0 & negligible (100\%): 0/3/0/0 & below (100\%): 3/0/0/0 \\
\quad of which writer quantile (9) & negligible (100\%): 0/5/0/4 & negligible (100\%): 0/6/0/3 & below (100\%): 9/0/0/0 \\
Pauses, fixed threshold (18) & negligible (100\%): 0/13/0/5 & depends (60\%): 0/6/9/3 & depends (81\%): 0/9/2/7 \\
\bottomrule
\end{tabular}
\end{table}

\subsection{The same diagnostic applied to other conclusions}
\label{sec:res-other}

The diagnostic applies to any conclusion drawn from the labels. We applied it to three more: a comparison of rates
between datasets, a count model within writers, and the decomposition of stall length.

\paragraph{Comparing rates between datasets.}
\label{sec:res-ratio}
Does a comparison of stall rates between datasets, such as between cohorts of a course, survive a change of
definition? Under the baseline the rates were
15.1 (KLiCKe), 13.1 (CS1 2019), 8.8 (CS1 2021) and 7.1 (CoAuthor) per hour; across the 89 definitions, the CS1 2019
rate alone ranged from 4.4 to 46.0. By the criterion set before the computations (25\% of definitions), the rank
order of the four datasets moved: it differed from the baseline order for 80 of 89 definitions, although the close
baseline rates of CS1 2019 and KLiCKe swap alone in 11 of them. We counted a ratio as reversed when its point value and its
interval lay on the other side of 1 from the baseline ratio, a description of each definition rather than a test over
the grid (Figure~\ref{fig:ratio}). For the two CS1 editions,
which share the course and the recording and which we had committed to report first, 20 of 89 definitions (22\%)
reversed the ratio, below the 25\% we had set after seeing point values (Section~\ref{sec:status}; 22 of 89 in
the time after the 100th interval, S2), but 23 (26\%), above it, when the first 60~s of each observation interval
were also excluded (S2); 18 of the 20 used a fixed $\theta$. Every fixed-$\theta$ definition reversed both comparisons of
CS1 with KLiCKe (66 and 53 reversals in all); CS1 2019 against CoAuthor reversed for 6 and the two
writing datasets for none (S3). Among the
fixed pause thresholds, the two CS1 editions reversed for none of 18, whereas the pairs with CoAuthor, whose pause rates
ranged from 42 to 498 per hour, reversed for 14 and 13 (S3), so by the same criterion the direction of a ratio changed in the pause family too. The
counts are not comparable across pairs, whose
baseline ratios and sample sizes differ.

\begin{figure}[t]
\centering
\begin{tikzpicture}
\begin{axis}[name=a, width=0.52\textwidth, height=5cm, xmin=0, xmax=91, ymin=-1.1, ymax=1.2, xticklabels={},
  ylabel={log rate ratio}, xlabel={definitions ordered by ratio}, grid=major, grid style={figband},
  title={(a) CS1 2019 / CS1 2021}, title style={font=\footnotesize}, tick label style={font=\footnotesize},
  label style={font=\footnotesize}]
  \addplot[only marks, mark=*, mark size=1pt, color=figc1, error bars/.cd, y dir=both, y explicit,
    error bar style={color=figc1!60}] table[x=i, y=lr, y error plus expr=\thisrow{hi}-\thisrow{lr},
    y error minus expr=\thisrow{lr}-\thisrow{lo}, discard if not={fam}{1}] {figdata/ratio_cs19_cs21.csv};
  \addplot[only marks, mark=square*, mark size=1pt, color=figc2, error bars/.cd, y dir=both, y explicit,
    error bar style={color=figc2!60}] table[x=i, y=lr, y error plus expr=\thisrow{hi}-\thisrow{lr},
    y error minus expr=\thisrow{lr}-\thisrow{lo}, discard if not={fam}{0}] {figdata/ratio_cs19_cs21.csv};
  \addplot[only marks, mark=star, mark size=3.5pt, thick, color=black] table[x=i, y=lr, discard if not={b0}{1}]
    {figdata/ratio_cs19_cs21.csv};
  \draw[black, thin] (axis cs:0,0) -- (axis cs:91,0);
\end{axis}
\begin{axis}[at={(a.east)}, anchor=west, xshift=0.9cm, width=0.52\textwidth, height=5cm, xmin=0, xmax=91,
  ymin=-1.1, ymax=1.2, xticklabels={}, xlabel={definitions ordered by ratio}, grid=major, grid style={figband},
  title={(b) CS1 2019 / KLiCKe}, title style={font=\footnotesize}, tick label style={font=\footnotesize},
  label style={font=\footnotesize}]
  \addplot[only marks, mark=*, mark size=1pt, color=figc1, error bars/.cd, y dir=both, y explicit,
    error bar style={color=figc1!60}] table[x=i, y=lr, y error plus expr=\thisrow{hi}-\thisrow{lr},
    y error minus expr=\thisrow{lr}-\thisrow{lo}, discard if not={fam}{1}] {figdata/ratio_cs19_kl.csv};
  \addplot[only marks, mark=square*, mark size=1pt, color=figc2, error bars/.cd, y dir=both, y explicit,
    error bar style={color=figc2!60}] table[x=i, y=lr, y error plus expr=\thisrow{hi}-\thisrow{lr},
    y error minus expr=\thisrow{lr}-\thisrow{lo}, discard if not={fam}{0}] {figdata/ratio_cs19_kl.csv};
  \addplot[only marks, mark=star, mark size=3.5pt, thick, color=black] table[x=i, y=lr, discard if not={b0}{1}]
    {figdata/ratio_cs19_kl.csv};
  \draw[black, thin] (axis cs:0,0) -- (axis cs:91,0);
\end{axis}
\end{tikzpicture}
\caption{Log ratios of stall rates with 95\% intervals. Circles: writer-quantile $\theta$; squares: fixed seconds;
star: baseline. A ratio reverses when its point value and its interval lie on the other side of zero from the baseline.}
\label{fig:ratio}
\end{figure}

\paragraph{Departure from a constant rate within writers.}
\label{sec:res-phi}
Do the counts of a writer's sub-units vary more than a constant rate per writer would make them vary? The answer
decides how much dispersion a count model of stalls has to allow. We measured the dispersion index $\varphi$ with the writer's own rate as the expectation (Table~\ref{tab:phi},
Figure~\ref{fig:phi}). Under the baseline, $\varphi$ exceeded 1 in all four CS1 cells (1.44 to 3.82) and not in
CoAuthor sessions (1.08). Across the 89 definitions, the median $\varphi$ was close to 1 in CS1 2019 blocks and CoAuthor
sessions (1.03 and 1.02), 1.21 in CS1 2021 blocks and 1.25 and 2.28 in the CS1 assignments; the interval exceeded 1 for
43 to 68 definitions in the CS1 cells and for 7 in CoAuthor. The baseline lay above the median of the 89 definitions in
every CS1 cell. The conclusion against 1.5 reversed, by the same counting rule, for 54 definitions in the CS1 2019
assignments, 49 in the CS1 2021 blocks, 19 (21\%) in the CS1 2021 assignments and none in CoAuthor; in CS1 2019 blocks
the baseline interval contained 1.5, so no reversal was counted. In the
CS1 2019 assignments and the CS1 2021 blocks, then, whether $\varphi$ lay above or below 1.5 depended on the definition;
in the other cells the reversals stayed under 25\% or could not be counted. Restricted to the time after the 100th
interval, the two cells were the CS1 2019 and CS1 2021 assignments (57 and 26 reversals), because the CS1 2021 blocks
could no longer be counted, and so they were when the first 60~s of each observation interval were also excluded
(57 and 24; S2). $\varphi$ compares variance with the mean
and does not show which assumption of a count model fails.

\begin{table}[t]
\centering\footnotesize
\caption{Dispersion index $\varphi$ within writers. Baseline with 95\% interval; median and range over the 89
definitions; number of definitions whose interval exceeds 1; number whose point value and interval lie on the other
side of 1.5 from the baseline (--: the baseline interval contains 1.5).}
\label{tab:phi}
\begin{tabular}{@{}l l r r r r@{}}
\toprule
Cell & Baseline & Median & Range & Interval $> 1$ & Reversed at 1.5 \\
\midrule
CS1 2019 blocks & 1.44 [1.35, 1.52] & 1.03 & 0.67--3.28 & 43 & -- \\
CS1 2019 assignments & 1.95 [1.75, 2.15] & 1.25 & 0.74--3.72 & 57 & 54 \\
CS1 2021 blocks & 1.82 [1.58, 2.11] & 1.21 & 0.56--4.66 & 46 & 49 \\
CS1 2021 assignments & 3.82 [3.13, 4.56] & 2.28 & 0.63--7.87 & 68 & 19 \\
CoAuthor sessions & 1.08 [0.98, 1.19] & 1.02 & 0.80--1.37 & 7 & 0 \\
\bottomrule
\end{tabular}
\end{table}

\begin{figure}[t]
\centering
\begin{tikzpicture}
\begin{axis}[width=\textwidth, height=5cm, xmin=0, xmax=461, ymin=0.4, ymax=8.2, ymode=log, log basis y=10,
  ytick={0.5,1,1.5,2,4,8}, yticklabels={0.5,1,1.5,2,4,8}, ylabel={$\varphi$}, xtick={45,137,229,321,413},
  xticklabels={CS1 2019 blocks, CS1 2019 assign., CS1 2021 blocks, CS1 2021 assign., CoAuthor sessions},
  tick label style={font=\scriptsize}, label style={font=\footnotesize}, grid=major, grid style={figband}]
  \addplot[only marks, mark=*, mark size=0.8pt, color=figc1, error bars/.cd, y dir=both, y explicit,
    error bar style={color=figc1!50}] table[x expr=(\thisrow{cell}-1)*92+\thisrow{i}, y=phi,
    y error plus expr=\thisrow{hi}-\thisrow{phi}, y error minus expr=\thisrow{phi}-\thisrow{lo}] {figdata/phi.csv};
  \addplot[only marks, mark=star, mark size=3.5pt, thick, color=black] table[x expr=(\thisrow{cell}-1)*92+\thisrow{i},
    y=phi, discard if not={b0}{1}] {figdata/phi.csv};
  \draw[black, thin] (axis cs:0,1) -- (axis cs:461,1);
  \draw[black, dashed] (axis cs:0,1.5) -- (axis cs:461,1.5);
\end{axis}
\end{tikzpicture}
\caption{The dispersion index within writers for the 90 definitions in each cell, ordered by value, with 95\%
intervals; star: baseline. Solid line: 1 (constant rate within a writer); dashed line: 1.5.}
\label{fig:phi}
\end{figure}

\paragraph{Components of stall length.}
\label{sec:res-length}
A map of what moves needs a side that stays. Is the decomposition of stall length as fragile as the rates? We fitted
the decomposition for all 90 definitions (Figure~\ref{fig:length}). The lengths are times from onset to resumption, so
a change of definition changes both the stalls and the point from which they are measured. Under the baseline, the stall-to-stall share was
0.93 (CS1 2019) and 0.92 (CS1 2021), and the writer share 0.03 in both, which can be read as the intraclass correlation
of writers. In both CS1 editions, the stall-to-stall share stayed the largest component in every one of the 89
definitions (0.88 to 1.00), at least 0.83 above the next component, and the writer share stayed small (0.00 to 0.05);
we did not compute intervals for these shares. The excluded unresumed stalls were 0.5--7\% of the CS1
stalls, depending on the definition. CoAuthor kept the same order in 86 of 89 definitions (S5). This is a weak contrast: most variance sitting at the lowest level is common for event-level
outcomes (S5).

\begin{figure}[t]
\centering
\begin{tikzpicture}
\begin{axis}[width=\textwidth, height=4.6cm, xmin=0, xmax=277, ymin=0, ymax=1.05, ylabel={share of variance},
  xtick={45,137,229}, xticklabels={CS1 2019, CS1 2021, CoAuthor}, tick label style={font=\footnotesize},
  label style={font=\footnotesize}, grid=major, grid style={figband},
  legend style={font=\scriptsize, at={(0.5,0.55)}, anchor=center}, legend cell align=left]
  \addplot[only marks, mark=*, mark size=0.9pt, color=figc1] table[x expr=(\thisrow{ds}-1)*92+\thisrow{i},
    y=stall] {figdata/lengths.csv};
  \addlegendentry{stall to stall}
  \addplot[only marks, mark=square*, mark size=0.9pt, color=figc2] table[x expr=(\thisrow{ds}-1)*92+\thisrow{i},
    y=writer] {figdata/lengths.csv};
  \addlegendentry{writer}
\end{axis}
\end{tikzpicture}
\caption{Shares of the variance of log stall length for the 90 definitions in each dataset, each series ordered by
value. The remaining shares (assignment, day, block or session) are in S5.}
\label{fig:length}
\end{figure}
